\textbf{Stochastic parametrisation in Lorenz-96.} The two-scale system was introduced by Lorenz~\cite{lorenz2006predictability}. Wilks examined the effects of stochastic parametrisations in it~\cite{wilks2005effects}, and Arnold et al.\ studied stochastic parametrisations and model uncertainty there~\cite{arnold2013stochastic}; our AR(1) closure follows this line, with the mean closure fitted by regression and the noise fitted to its residual. Outside this toy system, Shutts proposed a kinetic-energy backscatter algorithm for use in ensemble prediction systems~\cite{shutts2005kinetic}.

\textbf{Learned deterministic closures and coupled behaviour.} Rasp et al.\ learned all subgrid processes of a multiscale atmospheric model with a deep network~\cite{rasp2018deep}. Rasp observed that online simulations with offline-trained machine-learning parametrisations in Earth system models were frequently plagued by instabilities and biases, and proposed coupled online learning, a second training stage run alongside a high-resolution simulation, illustrating it in the Lorenz-96 model~\cite{rasp2019coupled}; a further approach trains parametrisations online for non-differentiable solvers with a neural emulator~\cite{frezat2023gradient}. We only train offline and measure what happens when the closure is coupled, so our stability result is limited to that setting.

\textbf{Learned stochastic closures.} Gagne et al.\ used generative adversarial networks as stochastic parametrisations in Lorenz-96~\cite{gagne2019machine}; Parthipan et al.\ combined a recurrent network with red noise in a probabilistic framework~\cite{parthipan2022using}. These are richer than our AR(1) closure and are out of scope here.

\textbf{Data-driven emulation of Lorenz-96.} Chattopadhyay et al.\ compared echo-state, feed-forward and LSTM networks for short-term prediction and long-term statistics of a multiscale Lorenz-96 system~\cite{chattopadhyay2020data}, and Brajard et al.\ combined data assimilation with a neural network to emulate a dynamical model from sparse noisy observations, with the 40-variable Lorenz-96 model as the test case~\cite{brajard2020combining}. Those works emulate the system itself; we instead keep the slow-variable equations and learn only the subgrid term.
